\documentclass[../main]{subfiles}
\begin{document}
\chapter{Unidades}
\img{images/00/presion.jpg}{Medición de presión de neumáticos-}
\info{
Contenido}{
Unidades de la hidrostática, Conversión de unidades, Métodos de conversión. En este trabajo práctico, se analizarán las unidades fundamentales de la hidroestática, como el bar, psi, etc. Se abordarán las definiciones de cada unidad, sus conversiones y aplicaciones en la práctica.}


\section{Introducción}

La hidroestática es la rama de la física que estudia los fluidos en reposo. Las unidades de medida son fundamentales para cualquier estudio científico, y en la hidroestática no es la excepción. 

En este trabajo práctico, se abordarán las unidades fundamentales de la hidroestática, como la presión, la densidad y el volumen. Se analizarán sus definiciones, conversiones y aplicaciones en la práctica.

\section{Unidades Fundamentales}

\subsection{Presión:}

La presión se define como la fuerza por unidad de área. La unidad fundamental de presión en el Sistema Internacional de Unidades (SI) es el pascal (Pa). 

Otras unidades de presión comunes son:

\begin{table}[h!]
    \centering
    \begin{tabular}{c|c|c}
* **Bar:** &1 bar = 100.000 Pa \\
* **Psi:** &1 psi = 6.894,76 Pa\\
* **Atmósfera:** &1 atm = 101.325 Pa \\
\end{tabular}
\end{table}
\subsection{Densidad:}

La densidad se define como la masa por unidad de volumen. La unidad fundamental de densidad en el SI es el kilogramo por metro cúbico (kg/m³). 

Otras unidades de densidad comunes son:
\begin{table}[h!]
    \centering
    \begin{tabular}{c|c|c}
    g/cm³ &1 g/cm³ & 1.000 kg/m³\\
lb/ft³ &1 lb/ft³ &16,02 kg/m³\\
    \end{tabular}
  
\end{table}



\subsection{Volumen:}

El volumen es la cantidad de espacio que ocupa un cuerpo. La unidad fundamental de volumen en el SI es el metro cúbico (m³). 

Otras unidades de volumen comunes son:

\begin{table}[h!]
    \centering
    \begin{tabular}{|c|c|c|} \hline
        Litro & 1 L & 1000cm³=0,001 m³ \\ \hline
        Galón & 1 gal & 0,00378541 m³\\ \hline
        Metro Cúbico & 1 m³ & 1000 L \\ \hline
    \end{tabular}
   \end{table} 
\section{Tabla general de Conversiones de Unidades}
\begin{center}
\begin{tabular}{|c|c|c|}
\hline 
\textbf{Unidad} & \textbf{Abreviatura} & \textbf{Equivalencia en SI} \\ \hline
Pascal & Pa & 1 Pa \\ \hline
Bar & bar & 1 bar = 100.000 Pa \\ \hline
Psi & psi & 1 psi = 6.894,76 Pa \\ \hline
Atmósfera & atm & 1 atm = 101.325 Pa \\ \hline
Kilogramo por metro cúbico & kg/m³ & 1 kg/m³ \\ \hline
Gramo por centímetro cúbico & g/cm³ & 1 g/cm³ = 1.000 kg/m³ \\ \hline
Libra por pie cúbico & lb/ft³ & 1 lb/ft³ = 16,02 kg/m³ \\ \hline
Metro cúbico & m³ & 1 m³ \\ \hline
Litro & L & 1 L = 0,001 m³ \\ \hline
Galón & gal & 1 gal = 0,00378541 m³ \\ \hline
\end{tabular}
\end{center}
\end{document}